\chapter{Réalisation et guide d'utilisation}
\label{ch:realisation}

\section{Stack technique déployée}

Sur \url{https://lexhub.space}, le navigateur parle à Caddy (HTTPS), qui route vers Next.js et FastAPI. FastAPI lit PostgreSQL, Redis et MinIO, délègue à Celery, et appelle ChatGPT, OCR, Whisper, Resend ou Turnstile selon le besoin.

\begin{table}[H]
\centering
\caption{Technologies LexHub}
\label{tab:stack}
\small
\renewcommand{\arraystretch}{1.25}
\begin{tabular}{|p{3.0cm}|p{10.2cm}|}
\hline
\rowcolor{tablehead}
\theadcell{Couche} & \theadcell{Choix déployé} \\
\hline
Langages & Python (API, Celery, masquage PII)~; TypeScript (Next.js)~; SQL \\
\hline
Interface & Next.js, Tailwind CSS, bilingue FR / AR \\
\hline
API & FastAPI, SQLAlchemy, Alembic \\
\hline
Données & PostgreSQL~16, Redis~7, MinIO \\
\hline
Tâches & Celery (OCR, e-mails, rappels) \\
\hline
Dictée & Web Speech (Chrome/Edge)~; Whisper (Firefox / audio) \\
\hline
IA texte & ChatGPT (OpenAI)~: correction et rédaction \\
\hline
OCR / export & OCR cloud~; Word (.docx) avec en-tête cabinet \\
\hline
Infra & Caddy, Docker Compose, VPS, GHCR (CI/CD) \\
\hline
Sécurité / mail & Turnstile (après 3 échecs)~; Resend (OTP, invitations, rappels) \\
\hline
\end{tabular}
\end{table}

\section{Préfixes API}
\label{sec:apis-prod}

Le frontend Next.js appelle FastAPI via ces préfixes (JWT).

\begin{table}[H]
\centering
\caption{Préfixes API LexHub}
\label{tab:apis}
\footnotesize
\renewcommand{\arraystretch}{1.22}
\begin{tabular}{|p{4.6cm}|p{8.6cm}|}
\hline
\rowcolor{tablehead}
\theadcell{Préfixe} & \theadcell{Rôle} \\
\hline
\texttt{/api/auth} & Inscription OTP, login, JWT, Turnstile, 2FA \\
\hline
\texttt{/api/dictation} & Whisper, correction, clarification, composition, révision \\
\hline
\texttt{/api/plaidoirie-}\newline\texttt{templates} & Brouillons, validation, assignation, export Word \\
\hline
\texttt{/api/documents} & Upload MinIO, OCR, questions sur pièce \\
\hline
\texttt{/api/cases} & Dossiers, analyse, clé de barre, planning lié \\
\hline
\texttt{/api/billing} & Plans, packs, sièges avocats +99~MAD, PAYG, factures \\
\hline
\texttt{/api/admin/billing} & Édition des plans (admin plateforme) \\
\hline
\texttt{/api/team} & Invitations e-mail, membres, rôles \\
\hline
\texttt{/api/workspace} & Tableau de bord cabinet \\
\hline
\texttt{/api/settings} & Charte, sécurité, suppression de compte \\
\hline
\texttt{/api/notifications} & Notifications in-app et digest \\
\hline
\texttt{/api/summarize} & Synthèses documentaires \\
\hline
\texttt{/api/corpus} & Corpus textuel cabinet \\
\hline
\texttt{/api/admin} & Cabinets, suspension, journal d'audit \\
\hline
\texttt{/api/ws/cabinet} & WebSocket temps réel (Redis) \\
\hline
\end{tabular}
\end{table}

\section{Rôles et espace cabinet}

Chaque rôle voit un périmètre distinct. Les détails des cas d'utilisation figurent au chapitre~\ref{ch:analyse}~; ici l'essentiel pour l'usage déployé.

\subsection{Avocat}
\begin{enumerate}[leftmargin=*,label=\arabic*.]
  \item Dicte, corrige, rédige, valide et exporte.
  \item Gère dossiers, pièces, planning et activité.
  \item Ne gère pas le paiement ni les invitations.
\end{enumerate}

\subsection{Admin cabinet}
\begin{enumerate}[leftmargin=*,label=\arabic*.]
  \item Dispose du même atelier que l'avocat.
  \item Invite l'équipe, gère le plan, le logo et le parrainage.
  \item Accède au menu Équipe et aux paramètres de facturation.
\end{enumerate}

\subsection{Assistant juridique}
\begin{enumerate}[leftmargin=*,label=\arabic*.]
  \item Consulte et télécharge les plaidoiries \textbf{validées} ou \textbf{finalisées} uniquement.
  \item Dépose des pièces et consulte le planning.
  \item N'utilise ni dictée, ni rédaction IA, ni synthèse, ni clé de barre.
\end{enumerate}

\subsection{Admin plateforme}
\begin{enumerate}[leftmargin=*,label=\arabic*.]
  \item Supervise LexHub depuis \texttt{/admin} sans lire le contenu des dossiers.
  \item Active, suspend ou purge un cabinet~; édite les plans~; consulte le journal d'audit.
  \item Les lifecycle emails partent automatiquement (Celery + Resend), sans action manuelle.
\end{enumerate}

\noindent Menu cabinet~: Plaidoirie, Dossiers, Documents, Planning, Activité, Équipe (admin cabinet). Paramètres~: cabinet, plan et usage, sécurité.

\section{Inscription et entrée en équipe}

L'inscription d'un \textbf{nouveau cabinet} : e-mail, OTP à six chiffres (Resend), nom du cabinet, mot de passe. Le cabinet est créé \textbf{actif} (plan Découverte, rôle ADMIN\_CABINET, essai). Aucun JWT n'est délivré à l'inscription ; l'utilisateur se connecte ensuite.

Pour rejoindre un cabinet \emph{existant}, l'admin envoie une \textbf{invitation par e-mail} (Resend). Le destinataire ouvre le lien, choisit son propre mot de passe, et accède au cabinet.

\section{Les cinq étapes de la plaidoirie}
\label{sec:guide}

Libellés du rail « Étape $X$ / $5$ » (capture en annexe, figure~\ref{fig:plaidoirie}).

\begin{enumerate}[leftmargin=1.4em]
  \item \textbf{Dossier \& pièces.} Dossier obligatoire. Fiche d'identité juridique. Ajout de pièces.
  \item \textbf{Dictée.} Web Speech (Chrome/Edge) ou Whisper (Firefox / fichier audio).
  \item \textbf{Transcription.} Texte brut éditable. Bouton \textbf{Corriger} (ChatGPT) : forme standard, sans résumer. Un texte sans phrase est \textbf{refusé}. Avant la rédaction : feuille \textbf{Clarification rapide}.
  \item \textbf{Rédaction IA.} Document éditable (ChatGPT), puis panneau de révision en langage naturel. L'avocat enregistre, valide ou assigne.
  \item \textbf{Export Word (.docx).} Document officiel : en-tête cabinet, typographie arabe. Après export réussi : statut \textit{finalized}.
\end{enumerate}

\section{Validation, assignation et assistant}

\begin{table}[H]
\centering
\caption{Collaboration autour d'une plaidoirie}
\label{tab:collab}
\small
\renewcommand{\arraystretch}{1.25}
\begin{tabular}{|p{3.0cm}|p{10.2cm}|}
\hline
\rowcolor{tablehead}
\theadcell{Action} & \theadcell{Effet} \\
\hline
Valider & Statut \textit{validated}~: tous les assistants du cabinet la voient. \\
\hline
Assigner & Valide si besoin et met la plaidoirie en avant pour l'assistant choisi. \\
\hline
Exporter Word & Statut \textit{finalized}. \\
\hline
Épingler & Mise en tête de liste. \\
\hline
Dévalider & Retour brouillon~: invisible des assistants. \\
\hline
\end{tabular}
\end{table}

Trois statuts~: brouillon, validée, finalisée. L'assistant ne voit que le validé ou le finalisé.

\section{Sécurité applicative}
\label{sec:securite}

\paragraph{Accès.}
Chaque requête authentifiée porte un JWT. Les routes sont filtrées par rôle. L'inscription d'un cabinet passe par un OTP e-mail (Resend). Les invitations d'équipe sont des liens nominatifs expirables. Une authentification à deux facteurs (TOTP) est optionnelle.

\paragraph{Anti-robot (Cloudflare Turnstile).}
Après \textbf{trois} échecs de login, le frontend charge le widget Turnstile ; l'API refuse le login tant que le jeton n'est pas validé auprès de Cloudflare.

\paragraph{Isolation.}
Toutes les entités métier sont rattachées à un \texttt{cabinet\_id}. L'API refuse une ressource hors locataire. Le navigateur ne parle qu'à FastAPI en HTTPS.

\paragraph{Pseudonymisation côté serveur (avant l'API).}
Coulisse backend, pas une étape de l'atelier~: juste avant l'appel à l'API OpenAI, \texttt{pii\_mask.py} remplace les identités par des codes (\texttt{PARTIE\_1}, \texttt{CIN\_1}, \ldots)~; \texttt{pii\_store.py} conserve la correspondance en PostgreSQL~; après la réponse, \texttt{unmask\_text} restaure les vrais noms pour l'UI et le Word. Le module traite noms (dossier + heuristiques dictée) et identifiants structurés (regex CIN / tél / e-mail). Ce n'est pas un NER entraîné~: une PII non reconnue peut partir en clair (marge d'erreur). Détail séquence~: diagramme~\ref{sec:seq-pii}.

\begin{figure}[H]
\centering
\includegraphics[width=0.88\textwidth]{fig_security_layers.png}
\caption{Couches de confidentialité LexHub en production (TLS, auth, RBAC, multi-tenant, PII, Adminer SSH).}
\label{fig:security-layers}
\end{figure}

\paragraph{Stockage des fichiers et de la base.}
Les pièces et logos sont dans MinIO. PostgreSQL reste sur le VPS, sans port ouvert sur Internet. Pour consulter la base en local, on utilise Adminer via un tunnel SSH (\texttt{http://127.0.0.1:8088}), jamais depuis l'URL publique du site.

\paragraph{Limitation d'abus et paramètres cabinet.}
Limitation de débit sur l'authentification et journal d'activité. L'admin plateforme peut suspendre un cabinet sans ouvrir les dossiers. Paramètres cabinet~: logo, 2FA, invitations, parrainage, digest e-mail, export planning.

\section{Lifecycle emails --- rétention et conversion}
\label{sec:emails}

Les e-mails automatiques de LexHub forment un levier \textbf{Growth / Retention}~: accompagner l'essai, rappeler les échéances, alerter sur les quotas et sécuriser le parcours (OTP, invitations). Canal unique~: \textbf{Resend}~; envoi via Celery. Le navigateur n'envoie jamais de mail directement. Le tableau~\ref{tab:emails} indique qui reçoit quoi et quand.

\begin{table}[H]
\centering
\caption{Lifecycle emails LexHub --- familles et déclencheurs}
\label{tab:emails}
\footnotesize
\renewcommand{\arraystretch}{1.2}
\begin{tabular}{|p{3.2cm}|p{2.8cm}|p{6.2cm}|}
\hline
\rowcolor{tablehead}
\theadcell{Famille} & \theadcell{Destinataire} & \theadcell{Quand} \\
\hline
OTP / sécurité & Utilisateur &
Inscription, reset mot de passe, changement d'e-mail. \\
\hline
Équipe & Admin cabinet et invité &
L'admin envoie l'invitation~; l'invité reçoit le lien~; l'admin est notifié s'il accepte ou refuse. \\
\hline
Plaidoirie & Assistant &
Validation, rappel agenda, digest. \\
\hline
Essai et abonnement & Admin cabinet &
Bienvenue essai, rappels fin d'essai, facture, alerte quota ou sièges. \\
\hline
Parrainage & Admin cabinet &
Lien de parrainage et récompense. \\
\hline
Impayé / renouvellement & Admin cabinet &
Échec de paiement, rappels, renouvellement. \\
\hline
\end{tabular}
\end{table}

\noindent Les messages portent un lien ou un code~; le corps du mail ne recopie pas le contenu des dossiers ni des plaidoiries.

